\section{Conclusions}
\label{sec:conclusions}
In this work, we propose improvements to a framework for pronunciation assessment that we recently introduced in \cite{cao24b_interspeech}.
%We propose two methods that allow the assessment of the pronunciation of a target speech sound without requiring knowledge of the corresponding speech segment boundaries.
We propose two methods with the goal of allowing the use of goodness of pronunciation (GOP) features extracted from modern high-performance automatic speech recognition models.
The main idea of our methods is to release the assumption that the GOP evaluation should be performed on a predetermined segment of speech corresponding to the target phone.
We do this either by adjusting the segment under test to the activations of the acoustic models (self-aligned GOP, GOP-SA) or by proposing a definition of GOP that is independent of the segmentation of test utterance (segmentation-free GOP, GOP-SF).

Our theoretical derivations make the underlying assumptions for GOP-SF explicit.
We also provide experimental results exploring a number of properties of our methods.
GOP-SA is always better than traditional GOP regardless of the characteristics of the acoustic model used for the evaluation.
GOP-SF obtains the overall best results for both the CMU Kids and for the speechocean762 material.
On the speechocean762 data we obtain state-of-the-art results keeping the MDD model constant.
We also show how the peakiness of the acoustic models affects the MDD results for the standard GOP definition and for the two proposed methods.
Finally we show that the performance of GOP-SF, that considers the full utterance to assess the pronunciation of the target phone, is not affected by the length of the context.

We believe that the proposed methods are potentially very appealing for phoneme-level pronunciation assessment, both because of their performance but also for their simple implementation and very low computational cost.
%In this work, we focus on providing theoretical explanation of the CTC-based framework of GOP introduced by~\cite{cao24b_interspeech}.  
%We hypothesize that the method GOP-CTC-AF is beneficial against alignment errors due to two reasons.
%The separation of information brought by CTC's peaky behavior while the alignment-free method considers all the alignments for GOP at the same time.  
%We provide experimental results by combining different models and methods to support our hypotheses. 
%In addition, we perform experiments to study the impact of GOP-AF regarding context length and activation length. 
%The results tell that GOP-CTC-AF is a measure which is robust to different context length and activation length while the performance of GOP-CE-AF can be optimized by the suggested normalization factor for activation length. 
%
%The limitations of GOP-AF for phoneme-level MDD derive from the limitation of the CTC where context information is not considered during training of the model.
%Since the context information usually has vital impact on MDD results as discussed earlier, investigating context dependent end-to-end ASR models such as RNN-T~\cite{graves2012sequence} to replace the CTC-based model is a direction for future study.
%Another potential direction is trying out more recent proposed SSL models, e.g., Hubert or WavLM~\cite{Chen2021WavLMLS} as backbone model and fine-tuning with a larger and noisy dataset (rather than the native English train-100-clean only here in this work). 
%Finally, since GOP is generally an ASR-based method in which the output from discriminative model is used. 
%This is a limitation of phoneme-level MDD because there could be other variabilities than phonemes to identify mispronounced speech, for example accent or dialect of the speakers.
%Therefore, using generative models for phoneme-level MDD is on our plan. 
